% Options for packages loaded elsewhere
\PassOptionsToPackage{unicode}{hyperref}
\PassOptionsToPackage{hyphens}{url}
%
\documentclass[
  a4paper,
]{article}
\usepackage{amsmath,amssymb}
\usepackage{lmodern}
\usepackage{iftex}
\ifPDFTeX
  \usepackage[T1]{fontenc}
  \usepackage[utf8]{inputenc}
  \usepackage{textcomp} % provide euro and other symbols
\else % if luatex or xetex
  \usepackage{unicode-math}
  \defaultfontfeatures{Scale=MatchLowercase}
  \defaultfontfeatures[\rmfamily]{Ligatures=TeX,Scale=1}
  \setmainfont[]{FreeSerif}
\fi
% Use upquote if available, for straight quotes in verbatim environments
\IfFileExists{upquote.sty}{\usepackage{upquote}}{}
\IfFileExists{microtype.sty}{% use microtype if available
  \usepackage[]{microtype}
  \UseMicrotypeSet[protrusion]{basicmath} % disable protrusion for tt fonts
}{}
\makeatletter
\@ifundefined{KOMAClassName}{% if non-KOMA class
  \IfFileExists{parskip.sty}{%
    \usepackage{parskip}
  }{% else
    \setlength{\parindent}{0pt}
    \setlength{\parskip}{6pt plus 2pt minus 1pt}}
}{% if KOMA class
  \KOMAoptions{parskip=half}}
\makeatother
\usepackage{xcolor}
\IfFileExists{xurl.sty}{\usepackage{xurl}}{} % add URL line breaks if available
\IfFileExists{bookmark.sty}{\usepackage{bookmark}}{\usepackage{hyperref}}
\hypersetup{
  pdftitle={N-srja},
  pdfauthor={Stefano Coretta},
  hidelinks,
  pdfcreator={LaTeX via pandoc}}
\urlstyle{same} % disable monospaced font for URLs
\usepackage[margin=1in]{geometry}
\usepackage{graphicx}
\makeatletter
\def\maxwidth{\ifdim\Gin@nat@width>\linewidth\linewidth\else\Gin@nat@width\fi}
\def\maxheight{\ifdim\Gin@nat@height>\textheight\textheight\else\Gin@nat@height\fi}
\makeatother
% Scale images if necessary, so that they will not overflow the page
% margins by default, and it is still possible to overwrite the defaults
% using explicit options in \includegraphics[width, height, ...]{}
\setkeys{Gin}{width=\maxwidth,height=\maxheight,keepaspectratio}
% Set default figure placement to htbp
\makeatletter
\def\fps@figure{htbp}
\makeatother
\setlength{\emergencystretch}{3em} % prevent overfull lines
\providecommand{\tightlist}{%
  \setlength{\itemsep}{0pt}\setlength{\parskip}{0pt}}
\setcounter{secnumdepth}{5}
\usepackage{expex}
\lingset{everyglpreamble=\bf}
\newcommand\SC[1]{\textsc{#1}}
\ifLuaTeX
  \usepackage{selnolig}  % disable illegal ligatures
\fi

\title{N-srja}
\author{Stefano Coretta}
\date{12/01/2020}

\begin{document}
\maketitle

\hypertarget{phonology}{%
\section{Phonology}\label{phonology}}

\textbf{Phonemes}. /t/, /s/, /r/, /j/, /h/.

\begin{itemize}
\tightlist
\item
  /t/: is {[}t{]} when non-syllabic, {[}n̩{]} when syllabic.
\item
  /s/: is {[}s{]} when non-syllabic, {[}s̩{]} when syllabic.
\item
  /r/: is {[}r{]} when non-syllabic, {[}r̩{]} when syllabic.
\item
  /j/: is {[}j{]} when non-syllabic, {[}i{]} when syllabic.
\item
  /h/: is {[}h{]} when non-syllabic, {[}ʰ{]} when preceded by a
  tautosyllabic consonant, {[}a{]} when syllabic.
\end{itemize}

\textbf{Syllabic structure}. C, CC, CCC.

\begin{itemize}
\tightlist
\item
  C: the consonant is syllabic. /t/ = {[}n̩{]}, /s/ = {[}s̩{]}.
\item
  CC: the second consonant is syllabic. /tt/ = {[}tn̩{]}, /jh/ =
  {[}ja{]}.
\item
  CCC: the third consonant is syllabic, unless two consecutive identical
  Cs. /tjh/ = {[}tja{]}, /rjt/ = {[}rjn̩{]}. When two consecutive
  identical Cs, the one that is second within the syllable is syllabic:
  /jtts/ = {[}itn̩s{]}, /trhh/ = {[}n̩rah{]}.
\end{itemize}

\textbf{Syllabification}. Consonants are assigned to syllables according
to the following algorithm:

\begin{enumerate}
\def\labelenumi{\arabic{enumi}.}
\tightlist
\item
  Assign the first C to its own syllable.
\item
  From the right, assign as many three-consonantal syllables as
  possible, unless this would create a second mono-consonantal syllable
  (*/C.C.CCC/ is not possible).
\item
  Assign the remaining consonants to two-consonantal syllables.
\end{enumerate}

Example:

\begin{enumerate}
\def\labelenumi{\arabic{enumi}.}
\tightlist
\item
  CCCCCC
\item
  (C)CCCCC
\item
  (C)CC(CCC)
\item
  (C)(CC)(CCC)
\end{enumerate}

/tsjhhr/ = {[}n̩.si.har{]}. /rrhtthj/ = {[}r̩.rʰn̩.tʰi{]}.

Roots are minimally disyllabic and maximally trisyllabic. Possible
roots:

\begin{itemize}
\tightlist
\item
  C.CC
\item
  C.CCC
\item
  C.CC.CC
\item
  C.CC.CCC
\item
  C.CCC.CCC (rare)
\end{itemize}

/tthj/ {[}n̩tʰi{]} `tree'. /tthjh/ {[}n̩taja{]} `trees'.

\end{document}
